% Related work -- paper plan Task 3.
% Source: research/02-differentiation.md §3 draft paragraph, with author-year
% stub entries replaced with refs.bib keys and the §"pre-generation decomposition
% of the query" wording corrected per ledger ruling R2 (the implementation
% decomposes the *draft*, not the query, before final response assembly).

\section{Related Work}
\label{sec:related}

\paragraph{Adaptive retrieval.}
Work on \emph{whether and when} to consult external evidence is well
established: Self-RAG trains reflection tokens to retrieve and self-critique
on demand \cite{asai2024selfrag}; CRAG gates generation on a
retrieval-quality evaluator with corrective search \cite{yan2024crag};
Adaptive-RAG routes queries by predicted complexity
\cite{jeong2024adaptiverag}; Rowen triggers retrieval on cross-response
inconsistency \cite{ding2025rowen}; and DRAGIN estimates the decoder's
real-time information need to retrieve mid-generation \cite{su2024dragin}.
None of these estimates \emph{per-claim hallucination risk}, and none judges
whether retrieved evidence is \emph{sufficient} for the specific claims
being written --- their triggers are relevance, uncertainty, complexity, or
information need.

\paragraph{Claim-level factuality.}
Atomic-claim decomposition has been formalized as \emph{evaluation}: FActScore
decomposes generations into atomic facts \cite{min2023factscore}, SAFE
verifies each fact with search-augmented agents
\cite{wei2024longformfactuality}, and HALoGEN adds high-precision atomic
verifiers and an error taxonomy \cite{ravichander2025halogen}. At inference
time, FACTOR allocates verification effort by claim-level risk
\cite{hassan2026factor} and SURE-RAG verifies claim-level evidence
sufficiency for selective answering \cite{qiu2026surerag}; CURE trains
claim-level confidence estimation \cite{liu2026cure} and Abbasli et al.
calibrate claim confidences post-hoc \cite{abbasli2026claimcalib}. These
systems evaluate, rank, or abstain --- they do not retrieve or repair.

\paragraph{Self-correction and localized repair.}
Self-correction ranges from full-response iteration
\cite{madaan2023selfrefine,shinn2023reflexion} to verification-then-full-
regeneration \cite{dhuliawala2023cove} and tool-grounded revision
\cite{gou2024critic}. RARR is the closest precedent for our repair stage:
it attributes and \emph{locally} edits unsupported content post-hoc while
preserving wording \cite{gao2023rarr}. DoctorRAG localizes failures in
agentic trajectories and repairs only the diagnosed point
\cite{jiao2026doctorrag}, OPT-Kraken selects repair granularity while
protecting correct content \cite{wu2026optkraken}, and DSVD rolls back and
re-decodes during generation \cite{guo2025dsvd}. RARR-style repair,
however, is strictly post-hoc: no risk estimate precedes it, no sufficiency
judgment gates it, and it emits no confidence estimate at all.

\paragraph{Structured claim--evidence state.}
Claim--evidence structures appear as an enterprise memory bank (ADORE
\cite{you2026adore}) and in fact verification over graphs
\cite{zhang2024causalwalk,guo2025muplon}; Enoki projects verified facts back
to spans \cite{rykov2026enoki}. In these systems the structure is a store or
an evaluation artifact, not a live control structure that simultaneously
decides retrieval, drives repair, and scores the emitted confidence.

\paragraph{Positioning.}
Each mechanism above has precedent, and we claim novelty only for the
integration: a single closed inference loop in which a shared claim state ---
draft-time decomposition, per-claim risk, sufficiency-gated retrieval, a
dynamic claim--evidence graph, localized healing, and per-claim
confidence --- governs every decision. The burden of proof is therefore
compositional: an ablation showing that each stage contributes and that the
integrated loop behaves differently from its parts (\S\ref{sec:results}).
